% !PRINTED 385
\weberpage{385}
%<<CONT
and denote the permutation thereby determined in a twofold way by

\weq{3}{\Pi_{1,2,3,8} = \Pi_{4,5,6,7},}

in that we stipulate that $\Pi_{\alpha_1\alpha_2\alpha_3\alpha_4}$ and $\Pi_{\beta_1\beta_2\beta_3\beta_4}$ shall mean the same thing when $\beta_1,\beta_2,\beta_3,\beta_4$ is any permutation of $\alpha_1,\alpha_2,\alpha_3,\alpha_4$, or when $\alpha_1,\alpha_2,\alpha_3,\alpha_4,\beta_1,\beta_2,\beta_3,\beta_4$ together comprise all eight figures.

By this symbol the permutation (2) is unambiguously designated, and the number of the distinct permutations of this kind amounts to exactly 35, so that we can represent the whole group $P$ by the subsidiary groups thus:

\weq{4}{P = S + \Sigma\, S\, \Pi_{\alpha_1\alpha_2\alpha_3\alpha_4},}

wherein $\alpha_1,\alpha_2,\alpha_3,\alpha_4$ run through all systems of four figures from the sequence $1,2,\dots,8$, in which an arbitrary figure, e.\,g. $\alpha_4 = 8$, can be held fixed.

To recognise the influence of $\Pi_{1,2,3,8}$ on any root $[i\,k]$, it suffices to consider the three roots $[1\,2], [1\,4], [4\,5]$, because $1,2,3$ on the one hand, $4,5,6,7$ on the other, occur entirely alike in $\Pi_{1,2,3,8}$.

These permutations are obtained simply from the remark that $[1\,2], [1\,4]$ are, by \S.\,99, the roots joined with $[2\,8]$ and $[4\,8]$ in the complex

\weq{5}{[2\,8][1\,2],\ [3\,8][1\,3],\ [4\,8][1\,4],\ [5\,8][1\,5],\ [6\,8][1\,6],\ [7\,8][1\,7]}

and that likewise $[4\,5]$ is the root joined with $[5\,8]$ in the complex

\weq{6}{[1\,8][1\,4],\ [2\,8][2\,4],\ [3\,8][3\,4],\ [5\,8][5\,4],\ [6\,8][6\,4],\ [7\,8][7\,4]}

But by the permutation (2) the complexes (5) and (6) pass over into the following, as one easily finds from the representation of the complexes in \S.\,99 [the complex (5) remains unchanged as a whole]:

\[
[3\,1][3\,8],\ [2\,1][2\,8],\ [4\,8][1\,4],\ [5\,8][5\,1],\ [6\,8][6\,1],\ [7\,8][7\,1]
\]

\[
[2\,3][1\,4],\ [3\,1][2\,4],\ [1\,2][3\,4],\ [5\,8][6\,7],\ [6\,8][5\,7],\ [7\,8][5\,6],
\]

from which one obtains the following permutations effected by $\Pi_{1,2,3,8}$:

\weq{7}{\begin{array}{l} [1\,2],\ [1\,4],\ [4\,5] \\ {}[3\,8],\ [1\,4],\ [6\,7]. \end{array}}

From (2) and (7) the following general and simple rule can now be derived.
